\documentclass{article}
\usepackage[utf8]{inputenc}
\usepackage[portuguese]{babel}
\usepackage{amsmath}
\usepackage{amssymb}
\usepackage{amsthm}
\usepackage{enumitem}
\usepackage{geometry}
\geometry{a4paper, margin=1in}

\title{Lista 3 - Sistemas Baseados em Conhecimento (MAC0444/MAC5778)}
\author{Leonardo Heidi Almeida Murakami \\ NUSP: 11260186}
\date{\today}

\begin{document}

\maketitle

\section*{1. Exercício 1}

Utilizando os conceitos atômicos \textit{Animal}, \textit{Peixe}, \textit{Peixívoro} e \textit{Gato}, os papéis \textit{temCauda}, \textit{temPata} e \textit{come} e a constante \textit{mimi}:

\subsection*{(a) Expressão dos conceitos em Lógica de Descrição}

Vamos traduzir cada conceito para a linguagem formal:

\begin{itemize}[leftmargin=*, itemsep=1em]
    \item \textbf{C1: Um animal que tem cauda.}
    \[
        \text{Animal} \sqcap \exists \text{temCauda}.\top
    \]
    
    \item \textbf{C2: Um animal que tem exatamente uma cauda e quatro patas.}
    \[
        \text{Animal} \sqcap ({=}\,1\text{ temCauda}.\top) \sqcap ({=}\,4\text{ temPata}.\top)
    \]
    
    \item \textbf{C3: Um animal que só come peixe.}
    \[
        \text{Animal} \sqcap \forall \text{come}.\text{Peixe}
    \]
    
    \item \textbf{C4: Um animal que só come coisas que só comem peixe.}
    \[
        \text{Animal} \sqcap \forall \text{come}.(\forall \text{come}.\text{Peixe})
    \]
\end{itemize}

\subsection*{(b) Expressão dos axiomas}

Agora vamos expressar cada axioma em Lógica de Descrição:

\begin{itemize}[leftmargin=*, itemsep=1em]
    \item \textbf{S1: Mimi é um gato que só come peixe.}
    \[
        \text{mimi} : \text{Gato} \sqcap \forall \text{come}.\text{Peixe}
    \]
    
    \item \textbf{S2: Gatos são animais que têm quatro patas.}
    \[
        \text{Gato} \sqsubseteq \text{Animal} \sqcap ({=}\,4\text{ temPata}.\top)
    \]
    
    \item \textbf{S3: Peixívoros são animais que só comem peixe.}
    \[
        \text{Peixívoro} \sqsubseteq \text{Animal} \sqcap \forall \text{come}.\text{Peixe}
    \]
\end{itemize}

\newpage

\subsection*{(c) S4 é consequência lógica de S1-S3?}

Queremos verificar se o axioma S4 pode ser deduzido dos axiomas anteriores:

\vspace{1em}

\noindent\textbf{S4: Mimi é um peixívoro de quatro patas.}
\[
    \text{mimi} : \text{Peixívoro} \sqcap ({=}\,4\text{ temPata}.\top)
\]

\vspace{1em}

\noindent\textbf{Prova:}

\vspace{0.5em}

Vamos construir um raciocínio passo a passo para mostrar que S4 realmente decorre de S1-S3.

\begin{enumerate}[leftmargin=*, itemsep=1em, label=\textbf{Passo \arabic*:}]
    \item O axioma S1 nos diz que:
    \[
        \text{mimi} : \text{Gato} \sqcap \forall \text{come}.\text{Peixe}
    \]
    
    Ou seja, podemos extrair duas informações sobre o mimi:
    \begin{itemize}[itemsep=0.3em]
        \item \(\text{mimi} : \text{Gato}\)
        \item \(\text{mimi} : \forall \text{come}.\text{Peixe}\)
    \end{itemize}
    
    \item Agora, sabendo que mimi é um gato e que o axioma S2 diz que todo gato é um animal com quatro patas, podemos concluir:
    \[
        \text{mimi} : \text{Animal} \sqcap ({=}\,4\text{ temPata}.\top)
    \]
    
    Em outras palavras:
    \begin{itemize}[itemsep=0.3em]
        \item \(\text{mimi} : \text{Animal}\)
        \item \(\text{mimi} : ({=}\,4\text{ temPata}.\top)\)
    \end{itemize}
    
    \item Juntando as informações dos passos 1 e 2, sabemos que mimi é um animal que só come peixe:
    \[
        \text{mimi} : \text{Animal} \sqcap \forall \text{come}.\text{Peixe}
    \]
    
    \item O axioma S3 define que:
    \[
        \text{Peixívoro} \sqsubseteq \text{Animal} \sqcap \forall \text{come}.\text{Peixe}
    \]
    
    Isso significa que ser peixívoro é exatamente ser um animal que só come peixe. Como mimi satisfaz essas condições (do Passo 3), podemos classificá-lo como peixívoro:
    \[
        \text{mimi} : \text{Peixívoro}
    \]
    
    \item Finalmente, combinando o fato de que mimi é um peixívoro (Passo 4) com o fato de que tem quatro patas (Passo 2):
    \[
        \text{mimi} : \text{Peixívoro} \sqcap ({=}\,4\text{ temPata}.\top)
    \]
\end{enumerate}

\vspace{1em}

\noindent\textbf{Conclusão:} S4 é consequência lógica de S1-S3.

\end{document}